% BHU PYQ -- Manually transcribed & error-corrected
% Paper: BCH-226 : Business Statistics
% Exam: B.Com. (Hons.) (Semester IV) (Special) Examination, 2014-15

\documentclass[11pt,a4paper]{article}
\usepackage[a4paper,top=2.5cm,bottom=2.5cm,left=2.8cm,right=2.8cm,headheight=14pt]{geometry}
\usepackage{amsmath,amssymb}
\usepackage[T1]{fontenc}
\usepackage[utf8]{inputenc}
\usepackage{lmodern,microtype,fancyhdr,enumitem,hyperref,lastpage,parskip}

\hypersetup{
    colorlinks=true,
    urlcolor=black,
    linkcolor=black,
    pdftitle={BCH-226: Business Statistics -- BHU B.Com. (Hons.) Sem IV 2014-15}
}

\pagestyle{fancy}
\fancyhf{}
\renewcommand{\headrulewidth}{0.4pt}
\renewcommand{\footrulewidth}{0.4pt}
\fancyhead[L]{\small   \textit{Banaras Hindu University}}
\fancyhead[R]{\small   \textit{BCH-226: Business Statistics}}
\fancyfoot[C]{\small Page  \thepage\ of \pageref{LastPage}}


\newlist{parts}{enumerate}{2}
\setlist[parts,1]{label=(\alph*),leftmargin=*,itemsep=5pt,topsep=3pt}
\setlist[parts,2]{label=(\roman*),leftmargin=*,itemsep=3pt,topsep=2pt}

\newcommand{\pts}[1]{\hfill   \textit{[#1]}}


\begin{document}


\begin{center}
  {\large\bfseries BANARAS HINDU UNIVERSITY}\\[2pt]
  {\normalsize B.Com. (Hons.) --- Semester IV (Special) Examination, 2014-15}\\[6pt]
  \rule{0.8\linewidth}{0.5pt}\\[6pt]
  {\large\bfseries Paper No. BCH-226: Business Statistics}\\[4pt]
  {\normalsize Time: 3 Hours\hfill Maximum Marks: 70}
\end{center}

\rule{\linewidth}{0.4pt}

\smallskip



\noindent   \textit{Instructions: Answer all questions. All questions carry equal marks (14 marks each).}

\bigskip




\noindent   \textbf{Question 1.}
Explain the procedure of constructing a cost of living index number. What are the difficulties associated with it?\pts{14}

\centerline{   \textbf{OR}}

Define Index Number and mention its uses and limitations.\pts{14}

\medskip\rule{\linewidth}{0.2pt}




\noindent   \textbf{Question 2.}
Differentiate between the Fixed Base and Chain Base methods of constructing index numbers and discuss their relative merits.\pts{14}

\centerline{   \textbf{OR}}

Calculate the real wage index numbers from the data given below. How much increase in money wages is required to maintain the real wages intact at the 2000 level?


\begin{center}
\small

\begin{tabular}{|l|c|c|c|c|c|c|c|}
\hline
   \textbf{Year} & 2000 & 2001 & 2002 & 2003 & 2004 & 2005 & 2006 \\
\hline
   \textbf{Wages (in Rs. 00)} & 20 & 24 & 35 & 36 & 36 & 38 & 40 \\
\hline
   \textbf{Price Index} & 100 & 150 & 200 & 220 & 230 & 250 & 250 \\
\hline
\end{tabular}
\end{center}\pts{14}

\medskip\rule{\linewidth}{0.2pt}

\pagebreak




\noindent   \textbf{Question 3.}
What are the components of a time series? How would you find out the trend values in a time series by the method of least squares?\pts{14}

\centerline{   \textbf{OR}}

Find out the value of the trend by three-yearly moving averages and plot them on a graph paper:


\begin{center}
\small

\begin{tabular}{|l|c|c|c|c|c|c|c|c|c|c|}
\hline
   \textbf{Year} & 1995 & 1996 & 1997 & 1998 & 1999 & 2000 & 2001 & 2002 & 2003 & 2004 \\
\hline
   \textbf{Value} & 10 & 15 & 12 & 18 & 15 & 22 & 19 & 24 & 20 & 26 \\
\hline
\end{tabular}
\end{center}\pts{14}

\medskip\rule{\linewidth}{0.2pt}




\noindent   \textbf{Question 4.}
What is seasonal variation in a time series? Describe the different methods you know to evaluate it and examine their relative merits.\pts{14}

\centerline{   \textbf{OR}}

Fit a parabolic curve of second degree ($y = a + bx + cx^2$) to the production data given below:


\begin{center}
\small

\begin{tabular}{|l|c|c|c|c|c|}
\hline
   \textbf{Year} & 2000 & 2001 & 2002 & 2003 & 2004 \\
\hline
   \textbf{Production (000 tonnes)} & 8 & 9 & 10 & 7 & 5 \\
\hline
\end{tabular}
\end{center}\pts{14}

\medskip\rule{\linewidth}{0.2pt}




\noindent   \textbf{Question 5.}
What do you understand by Statistical Quality Control? What are its advantages and disadvantages?\pts{14}

\centerline{   \textbf{OR}}

Write short notes on the following:

\begin{parts}
  \item Control Charts
  \item Process Control and Product Control
\end{parts}\pts{14}

\vfill
\rule{\linewidth}{0.4pt}

\begin{center}\small
\href{https://bhu-exam.vercel.app/}{Click here to visit our website}\\[4pt]
\href{https://me-aryan.vercel.app/}{Click here to visit our contributor}\\[4pt]
\href{https://quashan-me.vercel.app/}{Click here to visit our contributor}
\end{center}

\end{document}
